\section{The Eternal Race}

In \emph{The Eternal Race}, \textbf{Julius Evola} explains the doctrine of the ``\textbf{race of the spirit}``. Unlike the race of the soul, which is contingent and changeable, the race of the spirit is transhistorical. The way to access it is through the originary symbols and myths of a people. Evola describes the method, dangers, and value of this type of research.

Although the context is much different, this is not essentially much different from methods described by \textbf{Valentin Tomberg}; this undoubtedly arises from their deep and comprehensive knowledge of Hermetism. Remarkably, it is not unlike the ``phenomenology'' of \textbf{John Paul II} in the \emph{Theology of the Body}. The latter treats the first chapters of Genesis as the myth of the origins of the entire human race.

For Evola, however, the goal is to determine the spiritual roots of specific peoples. He expects a ``super race'' to arise, which is really Guenon's ``elite'' who, by recovering the original spiritual impulse of a people, will inaugurate a new civilization.

What follows is some background information for understanding the text.

\paragraph{Evolutionism}


Evola rejects ``evolutionism'' because it assumes that man evolved from a lower state to a higher state. This is a false metaphysical claim, although most people today assume it as a biological fact. What is too often called the ``theory of evolution'' is a misnomer; ``neo-Darwinism'' is a more accurate term (although there are alternatives), since it is based on the addition of Mendel's work on genes as the mechanism to explain heredity in Darwinism. However, strictly as science, there is no direction or aim to it. Humans can arise from apes just as readily as an ape from a human (and the latter is more observable.) There is only a problem when science oversteps its bounds.

\paragraph{Klages and Jung}


\textbf{Carl Jung} is easily misunderstood. In his autobiography, he explains that qua scientist, he can only make psychological claims. This is what Evola objects to. However, Jung also said that the archetypes have more than a psychological significance.

\textbf{Ludwig Klages} is a different story. \textbf{Lucian Tudor} describes Klages' philosophy of biocentrism in this way:

\begin{quotex}
Biocentrism posited an essential distinction between Seele (Soul?) and Geist (Spirit), which conflict each other in human life. The Geist is the nous, the pneuma, or the logos ... what we would call in ordinary speech the intellect, the reason, the spirit, the mind... . The Seele corresponds to the Greek psyche. It is the living principle, the vital spark ... and one with the body, soma. Biocentrism is a romanticist and anti-rationalist philosophy which poses the Soul as positive and the Spirit as negative. According to Biocentric theory, human beings had originally -- in primordial, ancient times -- lived ecstatically in accordance to the principle of Leben (Life, which is held as the ultimate value, to be distinguished from mere existence) and in a harmonious relation to the cosmos and the sensual world of images.

In Biocentric theory, at some point in history the force of the Spirit intruded into life and caused humans to use abstract, conceptual (as opposed to symbolic) thought and rational intellect, thus beginning the severing of body and Soul ... Biocentric philosophy also attacked Judeo-Christianity as a Logocentric religion opposed to Life and upheld ancient Paganism (which has a Dionysian character centered around vitalistic, feminine values) as the Biocentric religion of Life.
\end{quotex}


Obviously, Biocentrism is absurd from a Traditional point of view, so, for Evola, it is not only false but dangerous. However, it is implicit in much of neopagan reconstructions, which seeks to live in immediacy with no sense of the truly transcendent.

\paragraph{Positive Sources}


Evola mentions \textbf{Fustel de Coulanges} and \textbf{Bachofen} as being of fundamental importance. Evola's work will be understood better by those who have read those authors.

\paragraph{Text}


Let's now say something about research into the third grade of race, having the races of the spirit as its object. This is really the research of race at its ultimate root, wherever it is a question of normal civilizations and higher human lineages; the root, already connected to superpersonal, super-ethnic, metaphysical forces. For such research, the specific mode of conceiving both the sacred and the supernatural, as well as the relationship of man in respect to it, the vision of life in the highest sense, moreover, the entire world of symbols and myths constitute material as positive and objective as that of the racial features of the first degree such as facial features and cranial structures. Essentially in this area, the ``signs'' of ``vertical'', super-historical heredity are perceivable, which we said earlier; the special importance of this new research is confirmed from that point of view.

This, on the other hand, has more vast and precise possibilities of exploration of the origins and therefore of indicating the primary elements of races, than the other two degrees of race. The documents on which it is based can actually date back to the most distant prehistory, to that period which is precisely called ``mythical'' and which for this reason, is erroneously considered uncertain and deprived of importance by ``positive'' history. Furthermore, while the anthropological, archeological, and paleontological materials in themselves are mute, and while those of the research of the second degree are particularly subject to change, the myth and the symbol, in their atemporal and ahistorical nature, instead have an essential character of immutability, and therefore can often transmit elements conserving their original purity to a great extent. But naturally to reach that point, the entire trunk of prehistoric researches and whatever is connected to it, is to be sketched out by the new doctrine of race on a basis absolutely different from what applies to the older: on a pure, sacred, and no longer profane basis. It is necessary to proceed therefore to a complete revolution in the order of criteria and preconceptions that predominate in such a field, which, according to the usual trick, claim to be the measure of everything that may be considered ``serious'' and ``scientific''. And for the first thing, we repeat, it is necessary to get rid of the evolutionistic myth in all its forms, since it is evident that if one continues to believe that the further one goes back in time, the more one sinks into the horror of a bestial barbarity, the presumption of obtaining the points of reference valid for the present from the investigations of prehistory and of the mythical periods of origins, would be something insane. Wherever any evolutionistic premise is in use, to investigate the origins and to give prominence to the principle of heredity would lead fatally to aberrations like those of certain psychoanalytic exegetes such as Freud's \emph{Totem and Taboo}.

From the point of view here considered, we must say that the domain of the symbol and the myth, among us, is almost virgin territory. Giambattista Vico certainly did not create a school in Italy: rather, only secondary and often minor aspects of his theories have been adopted. Our official and unofficial culture, the one that calls itself ``serious'' and ``critical'' and that unfortunately is also largely represented in general education, considers symbol and myth to this day either as an arbitrary creation of pre-philosophical knowledge, or as something pertinent to inferior religious forms, or as a figurative and superstitious interpretation of merely natural phenomena, or, finally, as the formation of folklore---not to mention the discoveries of psychoanalysis and so called sociological schools, both typical creations of Judaism, that were introduced to us.

All these limitations and preconceptions must be overcome; otherwise, one must give up the fruits of the most fecund research in regard to race and primordial traditions. It is necessary to go back to conceiving myth and symbol as ancient, traditional man conceived it, i.e., as the characteristic expressions of a super-rational reality, objective in its own way, and almost as the sign, recognizable through every expert eye, of metaphysical forces that act in the depth of races, traditions, religions, and historical and prehistorical civilizations. Penetrating into the world of origins, while assuming such a point of view, is---we willingly concede---something not immune from dangers, because such a domain eludes the common means of control and criticism and, given the general lack of preparation in the field, every arbitrary and fantastic interpretation could have the right of citizenship in it. Germany has not failed to provide some examples of that. Without the armor of solid traditional principles and special qualification, quite different from that required for a ``critical'' research or for ``philosophical'' interpretation, move evil than good can come from the exploration into the question through deformations and corruptions.

In regard to general principles, if we want to profit from other people's experience in this aspect of research of the third degree, we must prevent an error, of which Klages and, in a certain measure, even Jung, can be considered the most significant proponents, which, by only recognizing the value of the symbol and myth as the object of a ``deep science'', it sees in them only a type of projection of the soul of the race, conceived irrationally as an expression of simple ``vital'' forces: ``Life'' (with the capital L) or the ``Collective Unconscious'' would be manifested in symbol and myth. That is false. And it is dangerous because it implies a romantic-naturalistic and rather one-sided conception of what race is, and must mean, for us. When it is a question of superior races, we repeat that the notion of race must be strictly joined to that of tradition, and in tradition, in its turn, the presence and the efficiency of meta-biological, metaphysical, not sub-rational but super-rational, forces must actually be recognized, acting formatively on purely physical and ``vital'' evidence and constituting the mystery of everything that, through race, assumes a fixed, unmistakable nature. Symbol and myth are signs of such deep forces of race, of which we just spoke, not of a type of irrational, instinctive and unconscious substrate of the ethnic group conceived in itself, a substrate that would be really make us think of the ``spirits'' or the \emph{totems} of primitive communities. In the face of confusions of this type, it is appropriate to recognize that some accusations against racial doctrine manifested as a type of new ``totemism'', a return to the spirit of the primordial horde, harmful for every true value of the personality, have a certain degree of justification.

Symbol and myth in our doctrine of race instead have the merit of documentary evidence through their capacity to lead us into the primary super-rational spiritual element of the races, to what is truly elemental in the world of origins. This element constitutes the leading thread for complementary investigations of various types. Custom, ethics, ancient law, and language certainly furnish other signs for the investigation of the third degree of race and for the racial interpretation of the history of civilization. But, even here, in order to obtain valid results it is necessary to remove the limitations of the modern mentality and to recognize that, in the ancient world, ethics, law, and customs were only chapters dependent on ``religion'': i.e., they reflect meanings and principles characteristic of a super-rational and sacred order. It is necessary to pick out the central point in this order, capable of conferring on everything else its own correct significance. When the investigation is closed to those forms taken in themselves, i.e., ethics, customs, law, language itself and art in the abstract, rather than as expressions above all of a given race of the body and the soul, then, by means of this, as applications or reflections of general meanings typical of tradition, the spiritual and animating force of race---when it is closed to that, again, it would remain in the are not of the original, but of the derived, not of the essential, but of the ancillary. In the face of so many of today's works, discouraged and losing themselves in the labyrinth of specialism and criticism without principles, the fundamental work of Fustel de Coulanges, like others of the same period and those of Bachofen, maintain, in this regard, in spite of all the imperfections dependent on the period in which they were written, a fundamental importance and are worth indicating the right direction for a series of studies that integrate them with the specific consideration of the race element.

Let us point out right now that to emphasize this original, spiritual element hidden in traditional myth and symbol, that, in races, goes beyond their merely biological, material and, fundamentally, even human, aspect, is very important even from the practical point of view. With it, actually, from what is conditioned by time and history and that, therefore, could only give rise to never experienced re-exhumations, almost ``commemorations'', one moves into the order of what, as essentially atemporal, is not to be considered of ``yesterday'', of a given ``history'' or ``prehistory'', but of a perennial relevance: to the eternal race. And it is exactly this race that can converted to an idée-force, capable of facilitating, through an awakening, by means of the law of attraction that attracts the similar to itself, the practical and creative tasks of applied doctrine of race, which is: the actualization of the ``super-race'' in the heart of the ``race'' as people, as the common type defined by a certain ethnic mixture, the re-emergence of the higher elements to the pure state and their formative reaffirmation, repeating the same mystery of the origins, in a new cycle of civilization.


\flrightit{Posted on 2015-05-24 by Cologero}

\begin{center}* * *\end{center}

\begin{footnotesize}
\begin{sffamily}
\texttt{ja on 2015-05-25 at 13:44 said:}

Klages philosophy is clearly derived from Nietzsche. ..a lot of neopaganism is just Nietzscheanism as opposed to a revival of Greco Roman thought.

\hfill

\texttt{Lucian Tudor on 2015-05-26 at 18:50 said:}

A more complete and effective exposition of Klages's thought can be found in the essays (available online) by Joseph Pryce and Lydia Baer. While I believe that authors like Ludwig Klages (or Jan Stachniuk, to cite a somewhat similar thinker) are useful for understanding Paganism as well as Lebensphilosophie, I think it is invalid to view their thought as the only representative form of Pagan philosophy. It is quite clear from Pryce's work, for example, that Stefan George and Klages had differing focal points for their Paganisms. Furthermore, I believe that more academic authors, historians of religion, such as Mircea Eliade provide the most true understanding of Paganism and its approach to the Sacred. However, it is best to compare and balance out a variety of viewpoints.

\hfill

\texttt{Cologero on 2015-05-27 at 07:38 said:}

Mr Tudor, we would certainly include Bachofen and Fustel de Coulanges as primary sources to understand the pagan mind and spirit. I disagree with your encyclopedic approach of comparing and balancing a variety of viewpoints. Either there is a standard of truth or there is not. Our primary interest, therefore, is whether a particular viewpoint adheres to or opposes that standard.

\hfill

\texttt{Lucian Tudor on 2015-05-27 at 08:57 said:}

Cologero, I definitely agree that there should be a standard of Paganism; a clear definition of what truly defines Paganism in all of its various forms across the world, which allows one to properly separate it from non-Pagan ideas and religions. What I meant by my remark is that to actually ascertain what Paganism truly is is not a simple task and can only be attained by reading a variety of authors' works, and attaining the ability to distinguish what is valid and invalid in each author after attaining an understanding of the essence of Paganism.

I believe that the more Biocentrist authors (Klages, Stachniuk, D.H. Lawrence, etc.) all grasped some truths about Paganism's defining characteristics, but Biocentrism misunderstands Paganism because it condemns all transcendence in favor of pure Life, whereas, by contrast, at the essence of Paganism is a complex, merged valuing of Life and Transcendence. Therefore, it can help to read authors like Klages to understand Paganism, but just reading these authors ultimately leads to an invalid understanding. Sure, Bachofen and Fustel de Coulanges should included in our list of authors that should be read in the study of Paganism, but recall that these are both academics from the 19th Century (and it is always proper to compare the works of different academics who have disagreements on the facts, as opposed to only reading one academic perspective on a subject) and they are also imperfect because different studies and facts have emerged since that time. It is necessary to include in one's study of Paganism modern intellecutals/scholars such as Eliade, Gilbert Durand, Nigel Pennick, Christopher Gérard, Alain de Benoist, etc. I cannot imagine a proper understanding of Paganism in the global sense being achieved without studying several different authors to some extent.

\hfill

\texttt{Mark Citadel on 2015-05-27 at 17:08 said:}

I am unsure as to how completely Evola rejects Evolution. Does he contend that the entire theory and its evidences are junk science, or is his claim more nuanced, positing his own historiography on the eternal race as a more purely metaphysical account rather than one rooted in what scientists would call `evidence'?

I'm leaning towards the former, that Evola considers the theory of evolution to be absolute junk, and affirms man created in present form.

\hfill

\texttt{Cologero on 2015-05-27 at 21:24 said:}

This is Evola's alternative to evolution: \textit{Esoteric Origin of the Species}\footnote{\url{http://www.gornahoor.net/?p=6261}}.


As we pointed out, the ``Theory of Evolution'' is a poor choice of words, since it implies a ``direction'' of evolution, from more primitive to more advanced. That is NOT what science qua science can show. Nevertheless, that is how it is viewed in the popular mind. Obviously, the Traditional viewpoint is that there is a degeneration over time, not a progression. Moreover, science has not proved, nor will it prove, a material cause for the human spirit.

\hfill

\texttt{mleow on 2015-05-30 at 17:36 said:}

Of interest i believe, the lost/last interview of Evola from french television has been published here:

\url{https://www.youtube.com/watch?v=QiCtdi5nCoA}


\hfill

\texttt{Cologero on 2015-05-31 at 12:45 said:}

Here is an interesting experiment to see if a contemporary political movement can be based on Traditional principles: \textit{The Esoteric Voice of Gabor Vona's Spiritual Mentor}\footnote{\url{http://budapesttimes.hu/2015/05/29/the-esoteric-voice-of-gabor-vonas-spiritual-mentor}}

\hfill

\end{sffamily}
\end{footnotesize}

